\section{Discussion and Limitations}
\label{sec:discussion-limitations}

A deployed system's limits are the part of it a practitioner most needs, and
stating them precisely is the same discipline that produced the results above,
applied to the boundary of what I can claim. Where a limit is a measurement I
have not made, I say so; where it is one I made and did not like, I report it.

\subsection{What the results support}

The central claim survives its strongest available test. Classifying by nesting
position rather than by label vocabulary yields recall of 1.000 at every
hierarchy level in all six states, including the two held out entirely, whose
documents print label words, code namespaces and level counts that appear
nowhere in the four the prompts were written against, and neither of which is
special-cased anywhere in the system. The two comparisons locate that result
rather than merely asserting it, and they locate it at the same boundary from
opposite sides: ablating the depth map degrades exactly the levels whose
identity is positional and leaves the others untouched, while the rule-based
extractor matches the language model at a document's top division and collapses
beneath it. Typography identifies that top division reliably and identifies
almost nothing below it.

\subsection{Corpus scope, and what the tiers do and do not mean}

\textbf{Every headline quality number in this paper is computed at the
\texttt{\_only\_subset} tier}: roughly one to two domains per state, 70 pages
in total, hand-trimmed and exhaustively annotated. This is favorable
preprocessing and I state it rather than leaving it to be inferred. It is the
tier that makes annotation affordable, and it is a genuine reduction in
coverage: a subset metric is not a whole-document metric, and no table at that
tier should be read as one. The one exception,
\S\ref{sec:experiments-scale-quality}, grades whole-document runs against a
golden of its own and is kept in its own table.

The scale measurements are a different tier and a different kind of claim. They
run at \texttt{\_trimmed}, which removes only non-standards matter (front
matter, guiding-principles essays, appendices, acknowledgements) and
\textbf{retains every standard in the document}. Kentucky's 52 processed pages
are all of Kentucky's standards, not 43\% of them, and Colorado's 41 are the
complete Ages 3--5 publication. A trimmed-to-published page ratio is therefore
never a coverage fraction, and I am explicit about this because the opposite
reading is the more natural one and would understate the result.

What remains limited about scale is \textbf{page count, as a cost and batching
question rather than a coverage one}. These runs chunk 18 and 17 times across
four \texttt{Map} iterations; their token and latency figures should not be
extrapolated to untrimmed page counts, and no document as large as Arizona's
217-page publication has been run end to end. Trimming also remains
preprocessing that a deployment would have to perform or absorb: the pipeline
was never asked to ignore 68 pages of Kentucky preamble on its own, and I do
not claim it would. Colorado carries one further qualification. Its
out-of-scope content is other \emph{age bands}, the birth-to-three and five-to-
eight portions of a 187-page publication, not standards removed by trimming. Every
Colorado figure in this paper is the Ages 3--5 band.

\subsection{What the goldens can and cannot establish}
\label{sec:discussion-goldens}

The detector goldens vary a great deal in density, and that governs what each
can establish. In the four development states they are deliberate spot checks:
5, 25, 7 and 8 annotated elements against detection sets of 67, 122, 61 and
33. Because the evaluation suite counts every unmatched in-scope detection as a
false positive, \emph{raw} precision on those states largely measures annotation
coverage rather than detector quality, and I do not report it as quality.
Precision instead comes from a manual audit of all 297 in-scope detections
across the six subset-tier documents, which found a single hallucinated
element.

The two held-out states are annotated densely rather than as spot checks, and
they separate two properties that are easily conflated. Nevada's golden is
content-exhaustive (46 elements covering every distinct structural element
the document prints) but not \emph{detection}-exhaustive, because the
document reprints several headings on a second page spread and the detector
correctly finds them twice; at 46 golden against 52 in-scope detections its raw
precision has a ceiling below 1 no matter how well it performs. Kentucky's
golden is the one that is detection-exhaustive, 44 against 44, so its precision
is a genuine hallucination rate rather than a coverage figure. It is also, for
that reason, what made quality-at-scale evidence available without new
annotation: those 44 elements are graded inside a trimmed-tier run in
\S\ref{sec:experiments-scale-quality}.

\paragraph{The trimmed-tier goldens support precision and not recall, and the
distinction is not a technicality.} The two Kentucky goldens used at that scale
(277 elements and 207 standards) were produced by seeding each row from a
recorded run and then checking it against the published PDF. Verifying emitted
rows establishes that what the system produced is correct, which is a precision
claim. It cannot establish recall, because it never searches for what the system
\emph{omitted}: an element no run emits is absent from the golden as well, and
invisible to any measurement made against it. Recall at that scale is therefore
reported against a verified reference set of 277 elements rather than against
the document, and the ceiling travels with the number wherever it appears.

Independent evidence bounds the gap at the middle levels without closing it.
Every distinct \texttt{Benchmark N.N} identifier the document prints (18 of 18)
and every \texttt{<Domain> Standard N} heading (15 of 15) appears in the golden,
recovered from the document text rather than from any run. The 202 indicators
carry no equivalent check, and that is where an omission would hide. Closing it
needs a cold read of the pages, reading a random sample without reference to
the golden and diffing, which would cost a fraction of annotating the document
and is the obvious next measurement.

\paragraph{Colorado deliberately carries only half of this.} Its golden remains
a seven-element spot check and is not being extended, so at the trimmed tier it
contributes recall and duplication evidence (zero elements absent, zero
duplicates, across four detection batches on the corpus's hardest layout) and
no precision figure at all. Its raw precision there is an annotation-coverage
artifact and is labeled as one in the recorded output. The corpus therefore has
exactly one state whose precision is a quality figure, at either tier, and that
is a real limit on how far the evidence generalizes.

\subsection{Nondeterminism, and why I do not claim determinism}

The pipeline is not deterministic, and the stability measurements should be read
as \textbf{lower bounds rather than estimates}. A defect firing in a minority of
runs survives five clean draws comfortably, so a zero in my stability table
means I did not observe instability at $n{=}5$, not that a component is stable.

Two of my own recorded results are the cautionary cases, and
\S\ref{sec:experiments-stability} reports both: a held-out code accuracy of
$46/46$ that five fresh runs give as 44, 45, 44, 44, 44, and a single-draw
attribution of Nevada's domain-code behavior to the Pass-1 sampler that $n{=}5$
per arm retired rather than confirmed. Both were the top of a distribution read
as its center. Where an experiment is repeated (the depth-map ablation, at
three runs per arm across six states) I report the range and note that the
direction never reverses, rather than quoting a single number.

The practical consequence is that correctness here rests on the validation
boundary rather than on the model agreeing with itself. The clearest instance is
California's intermittent structural-label defect: when it fires it produces
twelve malformed identifiers, all of which contain whitespace and are rejected
on shape before persistence.

\subsection{A known defect, measured and deliberately not repaired}
\label{sec:discussion-ageband-defect}

\S\ref{sec:experiments-scale-quality} reports that one field accounts for every
anomaly the trimmed tier introduces: the detector attributes a page banner to
individual elements, which defeats the deduplication key, which produces
duplicate standards, which the identifier resolver separates with fabricated
suffixes. I know the cause, I know the fix, and I have not applied it.

The fix is a prompt rule stating what the goldens already encode across all six
states: that an element's age band comes from a side-by-side column that
distinguishes it from its siblings, never from a banner or running header that
applies to a whole page. That is a general document-structure principle rather
than a per-state rule, so it belongs in the prompt by the discipline of
\S\ref{sec:method-llm-first}. Applying it changes a prompt, which changes the
hash that every recorded measurement in this paper is keyed to, and would
require re-recording the headline, ablation and comparison arms.

I report the defect rather than repair it, and state the reason plainly: the
cost of the repair is the invalidation of the evidence. That is the same
cost-gating that governs every other change to these two files, applied to a
finding of my own, and a reader is entitled to weigh it. Two things make the
deferral defensible rather than convenient. The defect does not affect any
subset-tier number, where the banner is not attributed at all. And it is
systematic rather than random: the count is 164 in each of three consecutive
runs, and 159 of those elements are common to all three, so its effect is
bounded and reproducible rather than a source of unquantified noise.

\subsection{Limitations of the measuring instruments}

Three are worth stating because each could mislead a reader who assumed
otherwise.

\textbf{The false-positive audit tests for invention.} Its verdicts turn on
whether a detected title appears in the source text, which is the right question
for a generative system and the wrong one for a copier: a rule-based
extractor cannot invent and therefore scores near 1.000 by construction. Applied
to the baseline the metric also detects something other than hallucination:
California's five flagged rows are column \emph{fusion}, where a table header
has been merged into an adjacent title, rather than fabrication. I report raw
precision for both arms instead, and treat the audit as informative only about
the language-model arm.

\textbf{The two evaluation suites do not grade the same detection.} The detector
suite re-runs detection in process against a frozen extraction (the direct
path) while the parser suite grades the deployed batched path's stored
output. On the recording this paper uses, Arizona's direct-path detection holds
67 elements where the batched output the parser suite reads holds 65; earlier
recordings have differed by more. The two suites also use
independently annotated goldens with different shapes, a flat element list and
nested normalized records. A detector number and a parser number from the same
results directory therefore describe related but not identical objects, and
should not be composed into a single end-to-end figure.

\textbf{Extraction quality bounds everything downstream.} The pipeline reads OCR
and PDF text-layer output; a heading the extractor fuses or drops is not
recoverable by any amount of downstream reasoning. Nevada's one confirmed
hallucination originates this way, in a multi-column table that flattens into
interleaved reading order.

\subsection{Scope}

The corpus is six US states of more than fifty jurisdictions, all in English,
all early learning standards. I make no claim about other document families,
other languages, or other countries' standards traditions, though the method
assumes only that a document has a nesting structure and prints some evidence
of it. Neither code nor corpus is released (\S\ref{sec:artifacts-statement}):
the source documents are third-party state-agency publications that are not ours
to redistribute, which limits direct reproducibility to the prompts and schema
disclosed in the appendices.
